\documentclass[12pt,a4paper]{article}

% --- Kodierung und Sprache ---
\usepackage[T1]{fontenc}
\usepackage[utf8]{inputenc}
\usepackage[ngerman]{babel}
\usepackage{lmodern}

% --- Mathematik ---
\usepackage{amsmath,amssymb,amsfonts,amsthm,mathtools}
\usepackage{bm}

% --- Layout ---
\usepackage[a4paper,margin=2.7cm]{geometry}
\usepackage{microtype}
\usepackage{setspace}
\onehalfspacing

% --- Tabellen / Links ---
\usepackage{booktabs}
\usepackage{hyperref}
\hypersetup{
	colorlinks=true,
	linkcolor=blue,
	citecolor=blue,
	urlcolor=blue,
	pdftitle={Von Hurwitz-Raeumen zu quantenartigen Hilbertraumstrukturen},
	pdfauthor={Thomas Hoffbauer},
	pdfsubject={Forschungsprogramm zur Bruecke zwischen Hurwitz-Operatorik und Quantenmechanik}
}

% --- Grafikeinbindung robust ---
\usepackage{graphicx}
\DeclareGraphicsExtensions{.pdf,.png,.jpg}
\graphicspath{{./}{fig/}{figs/}{images/}{img/}}
\newcommand{\includegraphicssafe}[2][]{%
	\IfFileExists{#2}{\includegraphics[#1]{#2}}{%
		\fbox{\parbox[c][3cm][c]{0.8\linewidth}{\centering Grafik \texttt{#2} nicht gefunden.}}%
	}%
}

% --- Theorem-Umgebungen ---
\theoremstyle{plain}
\newtheorem{satz}{Satz}[section]
\newtheorem{lemma}[satz]{Lemma}
\newtheorem{proposition}[satz]{Proposition}
\newtheorem{korollar}[satz]{Korollar}
\newtheorem{konjektur}[satz]{Konjektur}

\theoremstyle{definition}
\newtheorem{definition}[satz]{Definition}
\newtheorem{beispiel}[satz]{Beispiel}
\newtheorem{programm}[satz]{Programm}
\newtheorem{problem}[satz]{Problem}

\theoremstyle{remark}
\newtheorem{bemerkung}[satz]{Bemerkung}

% --- Makros ---
\newcommand{\OH}{\mathcal O_H}
\newcommand{\HN}{\mathcal H_N}
\newcommand{\Hinf}{\mathcal H_\infty}
\newcommand{\GN}{G_N}
\newcommand{\RN}{R_N}
\newcommand{\ZN}{\mathbb Z/N\mathbb Z}
\newcommand{\Zp}{\mathbb Z_p}
\newcommand{\Cp}{\mathbb C}
\newcommand{\Hp}{\mathbb H}
\newcommand{\elltwo}{\ell^2}
\newcommand{\nrd}{\operatorname{Nrd}}
\newcommand{\spec}{\operatorname{Spec}}
\newcommand{\im}{\operatorname{im}}
\newcommand{\Hom}{\operatorname{Hom}}
\newcommand{\tr}{\operatorname{tr}}
\newcommand{\id}{\mathrm{id}}

\title{\textbf{Von Hurwitz-Räumen zu quantenartigen Hilbertraumstrukturen}\\
	\large Ein Forschungsprogramm zwischen endlicher Operatorik, lokaler Arithmetik und Quantenmechanik}
\author{Thomas Hoffbauer}
\date{März 2026}

\begin{document}
	
	\maketitle
	
	\begin{abstract}
		Dieses Exposé formuliert ein präzises Forschungsprogramm zur Frage, ob und in welchem Sinn
		sich aus endlichen Quotienten der Hurwitz-Ordnung und den darauf definierten normgesteuerten
		Faltungsoperatoren eine Brücke zur Quantenmechanik auf Hilberträumen schlagen lässt.
		Ausgangspunkt ist eine bereits gesicherte endliche Spektralarchitektur: symmetrische
		Schrittmengen auf endlichen Gruppen erzeugen selbstadjungierte Faltungsoperatoren mit
		reellem Spektrum und spektraler Durchmischungsdynamik. Ziel ist es, die noch offenen
		Bausteine --- insbesondere die kanonische Definition der Normschalenoperatoren, ihre
		lokale arithmetische Interpretation, einen geeigneten Grenzübergang sowie eine
		quantenmechanische Zustands- und Messstruktur --- in eine kohärente Reihenfolge zu bringen.
		Der Text unterscheidet systematisch zwischen bewiesenem Kern, plausiblen Konjekturen
		und physikalischen Zusatzhypothesen.
	\end{abstract}
	
	\section{Einleitung}
	
	Die Quantenmechanik ist klassisch auf komplexen Hilberträumen formuliert. Zustände werden
	durch normierte Vektoren modulo Phase oder allgemeiner durch positive Spurklasseoperatoren
	beschrieben, Observablen durch selbstadjungierte Operatoren, und die Zeitentwicklung ist
	unitär. Diese Struktur ist außerordentlich erfolgreich, wirft aber weiterhin die Frage auf,
	ob sie tieferliegende diskrete, arithmetische oder geometrische Unterlagen besitzt.
	
	Parallel dazu tritt im hurwitzschen Kontext eine endliche Operatorarchitektur auf, die
	bemerkenswerte strukturelle Analogien besitzt: endliche Quotienten der Hurwitz-Ordnung,
	normgesteuerte Schrittmengen, selbstadjungierte Faltungsoperatoren und spektrale Dynamik.
	Diese Analogien reichen jedoch noch nicht aus, um eine physikalische Theorie zu begründen.
	Das Ziel dieses Exposés ist daher nicht, vorschnell eine neue Quantenmechanik zu behaupten,
	sondern den mathematischen Weg von einer endlichen hurwitzschen Spektralstruktur zu einer
	quantenartig interpretierbaren Hilbertraumstruktur in klaren Etappen zu formulieren.
	
	\section{Ausgangslage und Leitfrage}
	
	\subsection{Leitfrage}
	
	Die zentrale Frage lautet:
	
	\begin{quote}
		Kann aus einer kanonisch definierten Familie hurwitzscher Operatoren auf endlichen Quotienten
		ein Grenzobjekt konstruiert werden, das eine echte Hilbertraumstruktur mit selbstadjungierten
		Observablen, Wahrscheinlichkeitsinterpretation und reversibler Dynamik trägt?
	\end{quote}
	
	Diese Leitfrage zerfällt in drei Teilfragen:
	\begin{enumerate}
		\item Gibt es eine intrinsische, funktorielle und lokal verträgliche Definition der
		Normschalen \(\Sigma_{N,p}\subseteq G_N\)?
		\item Erzeugen die zugehörigen Operatoren \(T_{N,p}\) eine stabile Algebra mit
		kontrollierter lokaler und globaler Struktur?
		\item Lässt sich daraus ein quantenartiger Zustands- und Messformalismus gewinnen?
	\end{enumerate}
	
	\subsection{Methodische Vorsicht}
	
	Es ist entscheidend, zwischen drei Ebenen zu unterscheiden:
	\begin{enumerate}
		\item \textbf{Bewiesene endliche Operatorik:} selbstadjungierte Faltungsoperatoren,
		endliche Spektren, Spektrallücken, Durchmischung.
		\item \textbf{Arithmetisch-geometrische Konjekturen:} kanonische Normschalen, Hecke-Wirkung,
		Bruhat--Tits-Interpretation, Ramanujan-artige Schranken.
		\item \textbf{Physikalische Interpretation:} Zustände, Bornsche Regel, Dynamik,
		Tensorprodukt und Verschränkung.
	\end{enumerate}
	
	\section{Die gesicherte Grundarchitektur}
	
	\subsection{Hurwitz-Ordnung und endliche Quotienten}
	
	\begin{definition}[Hurwitz-Ordnung]
		Sei \(\OH\) die Hurwitz-Ordnung in der Hamiltonschen Quaternionenalgebra. Für \(N\geq 1\)
		definieren wir den endlichen Quotientenring
		\[
		R_N := \OH / N\OH.
		\]
		Aus \(R_N\) werde durch eine festgelegte Konstruktion eine endliche Gruppe \(G_N\) gewonnen.
	\end{definition}
	
	\begin{bemerkung}
		Die genaue Definition von \(G_N\) kann je nach Modellierung variieren. Für das folgende
		Programm genügt, dass \(G_N\) endlich ist und eine wohldefinierte Gruppenstruktur trägt.
	\end{bemerkung}
	
	\begin{satz}[Endlichkeit]
		Für jedes \(N\geq 1\) ist \(R_N\) ein endlicher Ring. Insbesondere ist \(G_N\) endlich.
	\end{satz}
	
	\begin{proof}[Beweisskizze]
		Da \(\OH\) als \(\mathbb Z\)-Modul Rang \(4\) besitzt, hat der Quotient \(R_N\) die Größe
		\(N^4\). Endlichkeit von \(G_N\) folgt aus seiner Konstruktion aus \(R_N\).
	\end{proof}
	
	\subsection{Hilbertraum auf endlicher Ebene}
	
	\begin{definition}
		Für jedes \(N\) sei
		\[
		\HN := \elltwo(G_N,\Cp)
		\]
		der komplexe Hilbertraum aller komplexwertigen Funktionen auf \(G_N\) mit dem Standardskalarprodukt
		\[
		\langle f,h\rangle := \sum_{g\in G_N} f(g)\,\overline{h(g)}.
		\]
	\end{definition}
	
	\begin{lemma}
		Jeder Raum \(\HN\) ist ein endlichdimensionaler komplexer Hilbertraum.
	\end{lemma}
	
	\begin{proof}
		Da \(G_N\) endlich ist, ist \(\elltwo(G_N,\Cp)\cong \Cp^{|G_N|}\).
	\end{proof}
	
	\subsection{Normschalen und Faltungsoperatoren}
	
	\begin{definition}[Normschale]
		Für eine Primzahl \(p\) sei
		\[
		X_p := \{x\in \OH : \nrd(x)=p\}
		\]
		die Normschale vom Niveau \(p\).
	\end{definition}
	
	\begin{definition}[Schrittmengen und Faltung]
		Ist \(\Sigma\subseteq G\) eine symmetrische Menge in einer endlichen Gruppe \(G\), so definieren wir
		\[
		(T_\Sigma f)(g):=\sum_{s\in\Sigma} f(gs).
		\]
		Der normierte Operator ist
		\[
		P_\Sigma := \frac{1}{|\Sigma|}T_\Sigma.
		\]
	\end{definition}
	
	\begin{lemma}[Selbstadjungiertheit]
		Ist \(\Sigma=\Sigma^{-1}\), so ist \(T_\Sigma\) selbstadjungiert auf \(\elltwo(G)\).
	\end{lemma}
	
	\begin{proof}
		Für \(f,h\in\elltwo(G)\) gilt
		\[
		\langle T_\Sigma f,h\rangle
		=
		\sum_{g\in G}\sum_{s\in\Sigma} f(gs)\overline{h(g)}
		=
		\sum_{g\in G} f(g)\overline{\sum_{s\in\Sigma} h(gs^{-1})}
		=
		\langle f,T_\Sigma h\rangle,
		\]
		weil \(\Sigma^{-1}=\Sigma\).
	\end{proof}
	
	\begin{korollar}
		Das Spektrum von \(T_\Sigma\) ist reell.
	\end{korollar}
	
	\begin{proposition}[Spektrallücke und Durchmischung]
		Ist \(\Delta(\Sigma)>0\) die Spektrallücke des normierten Operators \(P_\Sigma\), so gilt
		auf dem orthogonalen Komplement der konstanten Funktionen exponentielle Kontraktion:
		\[
		\|P_\Sigma^m f\|_2 \le (1-\Delta(\Sigma))^m \|f\|_2.
		\]
	\end{proposition}
	
	\section{Der eigentliche Engpass}
	
	\subsection{Das Kernproblem}
	
	Die gesicherte Theorie beginnt erst, \emph{wenn} eine wohldefinierte symmetrische Menge
	\(\Sigma_{N,p}\subseteq G_N\) gegeben ist. Das eigentliche offene Problem ist daher die
	Konstruktion solcher Mengen aus den globalen Normschalen \(X_p\).
	
	\begin{definition}[Kernproblem]
		Gesucht ist für \(p\nmid N\) eine kanonische Menge
		\[
		\Sigma_{N,p}\subseteq G_N,
		\]
		die zugleich folgende Eigenschaften besitzt:
		\begin{enumerate}
			\item Wohldefiniertheit modulo natürlicher Äquivalenzen,
			\item Symmetrie \(\Sigma_{N,p}=\Sigma_{N,p}^{-1}\),
			\item Ausschluss des neutralen Elements,
			\item gute Reduktion modulo \(N\),
			\item lokale Verträglichkeit für \(p\nmid N\),
			\item funktorisches Verhalten in \(N\) und \(p\).
		\end{enumerate}
	\end{definition}
	
	\begin{bemerkung}
		Gerade diese Definition ist im gegenwärtigen Stand die wichtigste Lücke. Solange sie nicht
		intrinsisch gelöst ist, bleibt die ganze Theorie zwar interessant, aber noch nicht natürlich.
	\end{bemerkung}
	
	\section{Das Forschungsprogramm}
	
	\begin{programm}[Von Hurwitz zu Hilbert]
		Wir formulieren das Gesamtziel als Kette von neun Stufen:
		\begin{enumerate}
			\item Konstruktion kanonischer Normschalen \(\Sigma_{N,p}\).
			\item Aufbau einer natürlichen Operatorfamilie \(T_{N,p}\).
			\item Nachweis lokaler arithmetischer Verträglichkeit.
			\item Vergleich mit Hecke- und Bruhat--Tits-Strukturen.
			\item Untersuchung von Kommutativität und gemeinsamer Spektralstruktur.
			\item Konstruktion eines Grenzhilbertraums \(\Hinf\).
			\item Definition von Zuständen und Observablen.
			\item Formulierung einer quantenartigen Dynamik.
			\item Aufbau einer Tensorstruktur für Mehrteilchensysteme.
		\end{enumerate}
	\end{programm}
	
	\subsection{Stufe I: Kanonische Normschalen}
	
	\begin{konjektur}[Kanonizität]
		Für jede ungerade Primzahl \(p\) und jedes \(N\) mit \(p\nmid N\) existiert eine intrinsische,
		symmetrische und funktorielle Menge \(\Sigma_{N,p}\subseteq G_N\), die aus den Norm-\(p\)-Elementen
		von \(\OH\) hervorgeht.
	\end{konjektur}
	
	\begin{bemerkung}
		Ein natürlicher Kandidat ist eine Doppelquotientierung
		\[
		\OH^\times \backslash X_p / \OH^\times
		\]
		mit anschließender Reduktion modulo \(N\). Hier stellt sich die Schlüsselfrage, ob diese
		Reduktion endlich-zu-eins mit kontrollierter Fasergröße ist.
	\end{bemerkung}
	
	\subsection{Stufe II: Die Operatorfamilie}
	
	\begin{definition}
		Ist \(\Sigma_{N,p}\) gegeben, so definieren wir
		\[
		T_{N,p}f(g):=\sum_{s\in \Sigma_{N,p}} f(gs)
		\qquad (f\in\HN).
		\]
	\end{definition}
	
	\begin{satz}
		Ist \(\Sigma_{N,p}\) symmetrisch, so ist \(T_{N,p}\) selbstadjungiert auf \(\HN\).
	\end{satz}
	
	\begin{bemerkung}
		Die Selbstadjungiertheit ist bereits auf endlicher Ebene gesichert. Offen ist jedoch die
		\emph{familienweise} Struktur der Operatoren.
	\end{bemerkung}
	
	\begin{konjektur}[Kommutativität]
		Für feste \(N\) kommutieren die Operatoren \(T_{N,p}\) und \(T_{N,q}\) für alle Primzahlen
		\(p,q\nmid N\), zumindest nach geeigneter Normierung oder Projektion.
	\end{konjektur}
	
	\begin{bemerkung}
		Ein positiver Nachweis würde den Übergang von bloßer diskreter Operatorik zu einer echten
		Observablenfamilie entscheidend stärken.
	\end{bemerkung}
	
	\subsection{Stufe III: Lokale Arithmetik}
	
	\begin{konjektur}[Lokale Split-Struktur]
		Für jede ungerade Primzahl \(p\) gilt lokal
		\[
		\OH \otimes \Zp \cong M_2(\Zp),
		\]
		wobei die reduzierte Norm zur Determinante wird. Unter dieser Identifikation gehen
		Norm-\(p\)-Klassen in Matrizen der Determinante \(p\) über.
	\end{konjektur}
	
	\begin{konjektur}[Lokale Nachbarschaft]
		Das Bild der globalen Norm-\(p\)-Klassen modulo lokalen Einheiten trifft genau eine
		kanonische Nachbarschaftsschale der Determinantenklasse \(p\).
	\end{konjektur}
	
	\begin{bemerkung}
		Dies wäre der erste echte Schritt in Richtung einer Bruhat--Tits-Interpretation. Er würde
		nicht sofort eine vollständige Gleichheit von Operatoren liefern, aber eine saubere lokale
		Geometrie bereitstellen.
	\end{bemerkung}
	
	\subsection{Stufe IV: Hecke-Interpretation}
	
	\begin{konjektur}[Hecke-Brücke]
		Die Operatoren \(T_{N,p}\) lassen sich nach geeigneter kanonischer Konstruktion mit endlichen
		Projektionen einer Hecke-Wirkung identifizieren.
	\end{konjektur}
	
	\begin{bemerkung}
		Diese Stufe ist mathematisch tief. Sie würde aus einem diskreten, arithmetisch motivierten
		Graphoperator einen echten automorphen oder lokal harmonischen Operator machen.
	\end{bemerkung}
	
	\subsection{Stufe V: Grenzübergang}
	
	\begin{problem}
		Die Familie \(\{\HN\}_N\) bildet noch keinen einzigen physikalischen Zustandsraum. Es muss
		ein Grenzübergang konstruiert werden.
	\end{problem}
	
	\begin{konjektur}[Grenzhilbertraum]
		Es existiert ein separabler komplexer Hilbertraum \(\Hinf\) sowie eine Familie kompatibler
		Einbettungen
		\[
		\iota_{N,M}:\HN \hookrightarrow \mathcal H_M
		\qquad (N\mid M),
		\]
		so dass die Operatoren \(T_{N,p}\) gegen selbstadjungierte Grenzoperatoren
		\[
		T_{\infty,p}:\Hinf \to \Hinf
		\]
		konvergieren.
	\end{konjektur}
	
	\begin{bemerkung}
		Hier sind mehrere Wege denkbar: direkter Limes, projektiver Limes, adelische Vervollständigung
		oder eine renormierte lokale Konstruktion. Welche dieser Varianten tragfähig ist, ist
		selbst Teil des Programms.
	\end{bemerkung}
	
	\subsection{Stufe VI: Zustände}
	
	\begin{definition}[Vorläufiger reiner Zustand]
		Ein vorläufiger reiner Zustand auf endlicher Ebene ist ein normierter Vektor
		\[
		\psi\in\HN, \qquad \|\psi\|=1.
		\]
	\end{definition}
	
	\begin{bemerkung}
		Dies ist zunächst nur eine \emph{formale} Analogie zur Quantenmechanik. Eine tiefere
		Frage lautet, ob der hurwitzsche Hintergrund eine genuin quaternionische Zustandsstruktur
		nahelegt oder ob der komplexe Hilbertraum die natürlichere Ebene bleibt.
	\end{bemerkung}
	
	\begin{problem}
		Soll die Theorie wirklich quaternionisch formuliert werden, so müssen Nichtkommutativität,
		Links-/Rechtsmodulstruktur und Phasenäquivalenz neu geklärt werden.
	\end{problem}
	
	\subsection{Stufe VII: Observablen und Bornsche Regel}
	
	\begin{definition}
		Eine vorläufige Observable ist ein selbstadjungierter Operator \(A\) in der von den
		\(T_{N,p}\) erzeugten Algebra.
	\end{definition}
	
	\begin{definition}[Born-ähnliche Regel]
		Ist \(A=\sum_\lambda \lambda P_\lambda\) die Spektralzerlegung von \(A\) und \(\psi\) ein
		normierter Zustand, so definieren wir
		\[
		\mathbb P_A(\lambda\mid \psi):=\langle \psi,P_\lambda\psi\rangle.
		\]
	\end{definition}
	
	\begin{bemerkung}
		Formal ist diese Regel auf endlicher Ebene immer möglich. Physikalisch interessant wird
		sie aber erst, wenn die Operatoren \(A\) nicht künstlich gewählt, sondern natürlich aus
		der hurwitzschen Struktur entstehen.
	\end{bemerkung}
	
	\subsection{Stufe VIII: Dynamik}
	
	\begin{definition}[Hamiltonoperator]
		Sei \(S\) eine endliche Menge von Primzahlen und \(a_p\in\mathbb R\). Dann definieren wir
		\[
		H_N := \sum_{p\in S} a_p T_{N,p}.
		\]
	\end{definition}
	
	\begin{satz}
		Ist \(H_N\) selbstadjungiert, so erzeugt
		\[
		U_N(t):=e^{-itH_N}
		\]
		eine unitäre Dynamik auf \(\HN\).
	\end{satz}
	
	\begin{proof}
		Dies ist eine Standardfolge des Spektralsatzes in endlichdimensionalen Hilberträumen.
	\end{proof}
	
	\begin{bemerkung}
		Hier entsteht erstmals eine echte quantenartige Zeitentwicklung. Ob sie mehr ist als ein
		formaler Transfer des Spektralsatzes, hängt davon ab, ob \(H_N\) arithmetisch kanonisch
		bestimmt werden kann.
	\end{bemerkung}
	
	\subsection{Stufe IX: Mehrteilchenstruktur}
	
	\begin{problem}
		Eine ernsthafte physikalische Theorie benötigt ein Tensorprodukt:
		\[
		\mathcal H_{AB} \cong \mathcal H_A \otimes \mathcal H_B.
		\]
		Im hurwitzschen Kontext ist keineswegs klar, wie eine solche Struktur natürlich entstehen
		soll.
	\end{problem}
	
	\begin{konjektur}
		Es gibt eine monoidale Struktur auf der Kategorie geeigneter hurwitzscher Operatorräume,
		die auf Hilbertraumebene ein natürliches Tensorprodukt induziert.
	\end{konjektur}
	
	\begin{bemerkung}
		Ohne diese Stufe bleibt das gesamte Programm bestenfalls eine Einteilchentheorie.
	\end{bemerkung}
	
	\section{Minimalziele und Testfälle}
	
	Bevor an eine große physikalische Deutung zu denken ist, sollten elementare Modelle
	rekonstruiert werden:
	
	\begin{enumerate}
		\item ein Zweiniveausystem,
		\item eine Spin-\(\frac12\)-artige Algebra,
		\item Interferenz in einem kleinen Invarianzunterraum,
		\item nichtkommutierende Observablen,
		\item diskrete Spektralaufspaltung,
		\item einfache Übergangsamplituden unter unitärer Entwicklung.
	\end{enumerate}
	
	\begin{bemerkung}
		Scheitert schon diese Ebene, so ist das Programm physikalisch zu schwach; gelingt sie,
		so lohnt der nächste Schritt.
	\end{bemerkung}
	
	\section{Was zu beweisen ist}
	
	Die mathematische Arbeit konzentriert sich auf die folgenden Hauptsätze.
	
	\begin{satz}[Zielsatz A: Kanonische Normschale]
		Für \(p\nmid N\) existiert eine kanonische symmetrische Menge \(\Sigma_{N,p}\subseteq G_N\),
		die funktoriell in \(N\) und lokal verträglich in \(p\) ist.
	\end{satz}
	
	\begin{satz}[Zielsatz B: Lokale Nachbarschaft]
		Die Mengen \(\Sigma_{N,p}\) entsprechen lokal einer kanonischen Nachbarschaftsschale der
		Determinantenklasse \(p\) in \(M_2(\Zp)\).
	\end{satz}
	
	\begin{satz}[Zielsatz C: Operatorfamilie]
		Die Operatoren \(T_{N,p}\) erzeugen für festes \(N\) eine strukturierte Algebra, idealerweise
		eine kommutative Hecke-artige Algebra.
	\end{satz}
	
	\begin{satz}[Zielsatz D: Hilbertraumlimes]
		Die Familie \((\HN,T_{N,p})_N\) besitzt einen wohldefinierten Grenzübergang zu einer
		selbstadjungierten Operatoralgebra auf einem separablen Hilbertraum \(\Hinf\).
	\end{satz}
	
	\begin{satz}[Zielsatz E: Quantenartige Interpretation]
		Für eine nichttriviale Klasse von Minimalmodellen trägt \(\Hinf\) eine Zustands-,
		Observablen- und Dynamikstruktur, die einer endlichdimensionalen Quantenmechanik
		äquivalent oder zumindest eng verwandt ist.
	\end{satz}
	
	\section{Was ausdrücklich noch nicht behauptet werden darf}
	
	Methodisch wichtig ist eine Negativliste. Gegenwärtig darf \emph{nicht} behauptet werden:
	
	\begin{enumerate}
		\item dass \(\Sigma_{N,p}\) bereits kanonisch konstruiert sei,
		\item dass \(T_{N,p}\) schon als echte Hecke-Operatoren identifiziert seien,
		\item dass die zugehörigen Graphen nachweislich aus Bruhat--Tits-Nachbarschaften stammen,
		\item dass die ganze Familie \(\{T_{N,p}\}_p\) kommutiert,
		\item dass bereits Ramanujan-artige Schranken vorliegen,
		\item dass damit schon eine physikalische Alternative zur Quantenmechanik erreicht sei.
	\end{enumerate}
	
	\begin{bemerkung}
		Gerade diese methodische Zurückhaltung ist kein Defizit, sondern wissenschaftliche
		Stärke. Sie trennt den bewiesenen Kern von den ambitionierten Brücken.
	\end{bemerkung}
	
	\section{Arbeitsplan}
	
	\subsection*{Phase 1: Arithmetische Konstruktion}
	\begin{enumerate}
		\item Explizite Berechnung von \(X_p\) für kleine Primzahlen.
		\item Untersuchung der Doppelquotientierung \(\OH^\times\backslash X_p/\OH^\times\).
		\item Kontrolle der Reduktion modulo \(N\).
		\item Test auf Symmetrie, Vielfachheiten und Neutralitätsausschluss.
	\end{enumerate}
	
	\subsection*{Phase 2: Numerische Operatorik}
	\begin{enumerate}
		\item Erzeugung von \(T_{N,p}\) für kleine \(N,p\).
		\item Berechnung von Spektren und Kommutatoren.
		\item Suche nach gemeinsamen Eigenbasen.
		\item Untersuchung spektraler Lücken und Multiplizitäten.
	\end{enumerate}
	
	\subsection*{Phase 3: Lokale Struktur}
	\begin{enumerate}
		\item Konstruktion expliziter Isomorphismen \(\OH\otimes\Zp \cong M_2(\Zp)\) für kleine \(p\).
		\item Vergleich der globalen Normklassen mit lokalen Determinantenbahnen.
		\item Test auf Bruhat--Tits-artige Nachbarschaften.
	\end{enumerate}
	
	\subsection*{Phase 4: Hilbertraum und Dynamik}
	\begin{enumerate}
		\item Entwicklung eines Limes \(\Hinf\).
		\item Konstruktion natürlicher Hamiltonoperatoren \(H_N\).
		\item Analyse unitärer Dynamik auf kleinen invarianten Unterräumen.
	\end{enumerate}
	
	\subsection*{Phase 5: Physikalische Minimalmodelle}
	\begin{enumerate}
		\item Zweiniveausysteme.
		\item Spin-\(\frac12\)-Modelle.
		\item Interferenz und Übergangswahrscheinlichkeiten.
		\item erste Tensorproduktansätze.
	\end{enumerate}
	
	\section{Schluss}
	
	Die Idee einer Abbildung zwischen Quantenmechanik im Hilbertraum und einem hurwitzschen
	Raum ist mathematisch ernstzunehmen, aber nur unter einer klaren Bedingung:
	Sie darf nicht als fertige Identifikation präsentiert werden, sondern muss als
	mehrstufiges Programm verstanden werden.
	
	Der gegenwärtig gesicherte Kern ist eine endliche Spektraltheorie symmetrischer
	Faltungsoperatoren auf Quotienten der Hurwitz-Ordnung. Der entscheidende nächste
	Schritt besteht in der intrinsischen Konstruktion der Normschalen \(\Sigma_{N,p}\).
	Erst danach werden die Brücken zu lokaler Arithmetik, Hecke-Theorie, Hilbertraumlimes
	und quantenartiger Interpretation sinnvoll und prüfbar.
	
	In diesem Sinn ist das Programm weder bloße Metapher noch schon vollendete Theorie,
	sondern ein präzise formulierbares mathematisches Forschungsprojekt.
	
\end{document}